% Generated draft: requires mathematical and source review.
\documentclass[11pt]{article}
\usepackage[margin=25mm]{geometry}
\usepackage{fontspec}
\setmainfont{Latin Modern Roman}
\usepackage{amsmath,amssymb,mathtools,unicode-math}
\setmathfont{Latin Modern Math}
\usepackage{enumitem,needspace,xurl,hyperref}
\setlength{\parindent}{0pt}
\setlength{\parskip}{5pt}
\setlength{\emergencystretch}{3em}
\newcommand{\N}{\mathbb{N}}
\newcommand{\Z}{\mathbb{Z}}
\newcommand{\Q}{\mathbb{Q}}
\newcommand{\R}{\mathbb{R}}
\newcommand{\C}{\mathbb{C}}
\newcommand{\F}{\mathbb{F}}
\newcommand{\A}{\mathbb{A}}
\newcommand{\PP}{\mathbb{P}}
\newcommand{\E}{\mathbb{E}}
\newcommand{\eps}{\varepsilon}
\newcommand{\dd}{\,\mathrm{d}}
\begin{document}
\section*{UnsolvedMath ID: 20000455}
\textbf{Source problem code:} catalogue entry
\subsection{Definitions and mathematical statement}
A subset \(D\subseteq k\) is Diophantine over a number field \(k\) if there are finitely many polynomials \(P_j\in k[X,Y_1,\ldots,Y_m]\) such that
\[
 x\in D\ \Longleftrightarrow\ \exists y\in k^m\quad P_j(x,y)=0\text{ for every }j.
\]
Equivalently it is a coordinate projection of the \(k\)-points of an affine algebraic set. For a fixed archimedean place \(v\) put \(B_v=\{x\in k:|x|_v\le1\}\). Shlapentokh asks, for the source's nonreal number-field setting, to construct such a set \(S\subseteq k^n\) whose projection is exactly \(B_v\) for some chosen archimedean absolute value. The place is fixed for the whole ball, rather than chosen separately for each element. The extracted condition ``\(k\ 6\ \subset\ \R\)'' is read as a damaged non-containment symbol; whether this refers to a specified nonreal embedding or to an intrinsic condition on \(k\) is not spelled out. Likewise ``\(|k|\le1\)'' is understood as the bound on the element being tested, not on the field itself.

The supplied partial result fixes an embedding \(\sigma:k\hookrightarrow\C\) with nonreal image stable under complex conjugation. Let \(c\) be its induced involution and \(F=k^c\), embedded in \(\R\) by \(\tau=\sigma|_F\). It asserts
\[
 \{x\in k:|\sigma(x)|\le1\}\text{ is Diophantine over }k
 \quad\Longleftrightarrow\quad F\text{ is Diophantine over }k.
\]
For a nonzero \(\theta\in k\) with \(c(\theta)=-\theta\), the reported forward construction recovers equality of complex moduli from the ball, with a separate treatment of the zero-denominator case, and then uses
\[
 F=\{x\in k:|\sigma(x+\theta)|=|\sigma(x-\theta)|\}.
\]
For the reverse direction, write \(x=a+b\theta\), \(a,b\in F\), and set \(d=-\theta^2\in F\), so \(\tau(d)>0\). For the ordinary complex modulus the norm identity is
\[
 |\sigma(a+b\theta)|^2=\tau(a^2+d b^2).
\]
The report invokes Diophantine definability of the ordering at this chosen real embedding of \(F\) to reconstruct the ball when \(F\) is Diophantine over \(k\). Its summary writes a variant \(w(a^2+d b^2)\) without defining \(w\); the displayed standard-modulus identity here specifies the embedding rather than assigning an unsupported meaning to \(w\). The full definability constructions and their primary references are absent.

The report explicitly leaves the general intended complex-place construction open, in particular without conjugation stability or a Diophantine fixed real subfield. The source remarks that the problem would have consequences for Hilbert's tenth problem over rings of integers and is probably hard because of a near-equivalence, but does not state a precise reduction or equivalence theorem.
\subsection{Short English statement}
Define an exact archimedean unit ball over a nonreal number field by polynomial equations and existential projection. The supplied report reduces the conjugation-stable complex-embedding case to Diophantine definability of its fixed real subfield, while retaining the broader construction as open.
\subsection{Sources}
\begin{itemize}
\item Local UnsolvedMath registry, permanent ID \(20000455\), source code \texttt{catalogue entry}; AIM algebraic number theory corpus, \texttt{aim-algebraic-number-theory-notes.json}, notes[106]. Attribution: Ulam AI / UnsolvedMath Contributors, CC BY 4.0.
\item AIM \emph{Rational and integral points on higher dimensional varieties}, Question 22, attributed to Shlapentokh: \url{https://aimath.org/WWN/qptsurface2/qptsurface2.pdf}. Its remark concerns Hilbert's tenth problem and ends with ``Problem/''. The supplied verification refers to official PDF/HTML wording but ends before quoting it; this draft claims no new source inspection.
\item Original literature field is empty; no authoritative bibliographic reference for the fixed-subfield equivalence or chosen-order definability is supplied.
\item Registry dated August 14, 2026: statement exact, status \texttt{partially\_solved}, classification \texttt{PARTIAL-PROGRESS}; attempt 1 records problem status open, partial result, rigor proved, novelty confidence low. The candidate equivalence is a machine-generated research aid, not peer reviewed, requiring independent verification before citation.
\end{itemize}
\end{document}
